\section*{Bab \ref{ch:b3:sensory-systems} --- Sistem Indra}

\begin{solution}{exo:b3:sensory-systems:1}
Lewat kawatnya, bukan pesannya: tiap saraf indra adalah sebuah \omterm{def:b3:sensory-systems:principles}{jalur berlabel}
yang berakhir di daerahnya sendiri di dalam otaknya, dan apa pun yang tiba lewat
saraf optiknya dibaca sebagai cahaya --- dan itulah sebabnya tekanan pada bola mata
menghasilkan sebuah kilatan dan sebuah pukulan pada telinga menghasilkan sebuah dengingan.
\end{solution}

\begin{solution}{exo:b3:sensory-systems:2}
Di dalam gelap, GMP siklik menahan saluran kation tetap terbuka dan sebuah arus masuk
yang tetap, yakni ``arus gelap'', menjaga selnya tetap terdepolarisasi dan melepaskan glutamat.
Cahaya mengaktifkan \omterm{def:b3:sensory-systems:retina}{rodopsin}, transdusin, dan fosfodiesterase, yang
memusnahkan cGMP; salurannya menutup, arus masuknya berhenti, dan
selnya terhiperpolarisasi lalu melepaskan lebih sedikit --- cahaya diisyaratkan oleh sebuah
penurunan.
\end{solution}

\begin{solution}{exo:b3:sensory-systems:3}
$I = I_{0}10^{L/10}$: \qty{20}{dB}, \qty{1e-10}{W/m^2}; \qty{60}{dB},
\qty{1e-6}{W/m^2}; \qty{110}{dB}, \qty{0.1}{W/m^2}. Ternyaring ke terlembut:
$10^{9}$.
\end{solution}

\begin{solution}{exo:b3:sensory-systems:4}
Manis, gurih, dan pahit lewat reseptor tergandeng-protein~G (satu
manis, satu gurih, sekitar dua puluh lima pahit); asin lewat saluran natrium
ENaC; asam lewat sebuah saluran proton. Kebanyakan ``cita rasa'' adalah
bau molekul mudah menguap yang mencapai hidungnya dari mulutnya; dengan
hidungnya tersumbat hanya kelima rasanya yang tersisa dan makanannya terasa hambar.
\end{solution}

\begin{solution}{exo:b3:sensory-systems:5}
Fechner: $S = (1/0.02)\ln(10\,000/100) = 50\ln 100 \approx 230$ langkah.
Mencacah: $1.02^{n} = 100$ memberi $n = \ln 100/\ln 1.02 = 233$.
\end{solution}

\begin{solution}{exo:b3:sensory-systems:6}
$P = 1 - e^{-a}\sum_{k=0}^{5}a^{k}/k!$: $a = 2$: $0.017$; $a = 6$:
$0.55$; $a = 12$: $0.98$. Separuh kilatannya terlihat pada $a \approx 5.7$;
dengan penyerapan \qty{10}{\%} kilatannya harus menghantarkan sekitar $57$
foton pada korneanya.
\end{solution}

\begin{solution}{exo:b3:sensory-systems:7}
$55/3\times 1.3 \approx 24$; dalam \omterm{def:b3:sensory-systems:cochlea}{desibel} tekanan, $20\log_{10}24
\approx \qty{28}{dB}$. Tulang yang berlekatan tidak dapat meneruskan getarannya, sehingga
bunyinya mencapai \omterm{def:b3:sensory-systems:cochlea}{koklea} hanya lewat hantaran menembus tengkoraknya,
tanpa bati telinga tengahnya: yakni sebuah kehilangan hantaran sekitar
\qtyrange{30}{50}{dB}.
\end{solution}

\begin{solution}{exo:b3:sensory-systems:8}
Saluran \omterm{def:b3:sensory-systems:cochlea}{sel rambutnya} dibuka langsung oleh \omterm{def:b3:sensory-systems:cochlea}{kaitan ujung} yang menarik
padanya --- tanpa pesuruh, tanpa enzim --- sehingga ia menanggapi dalam hitungan mikrodetik;
fotoreseptornya melipatgandakan satu foton lewat dua tahap enzimatis, yang
memakan puluhan milidetik. Pendengarannya memperoleh kecermatan waktu yang diperlukan untuk
menempatkan bunyi lewat penundaan antartelinga dan untuk mengikuti fasenya; penglihatannya memperoleh
kepekaan untuk mencacah foton tunggal, dan membayarnya dengan kecepatan.
\end{solution}

\begin{solution}{exo:b3:sensory-systems:9}
$H(0.05) = -0.05\ln 0.05 - 0.95\ln 0.95 = 0.150 + 0.049 = 0.199$.
Manusia: $\ln\binom{400}{20} \approx 400\times 0.199 = 80$, yakni sekitar
$10^{34}$ pola. Mencit: $\ln\binom{1100}{55} \approx 1100\times 0.199
= 219$, yakni sekitar $10^{95}$.
\end{solution}

\begin{solution}{exo:b3:sensory-systems:10}
Rata-ratanya $500\times 0.1/40 = 1.25$ peristiwa spontan tiap jendela. $P(k
\le 5) = e^{-1.25}(1 + 1.25 + 0.78 + 0.33 + 0.10 + 0.03) = 0.998$, sehingga
$P(k \ge 6) \approx 2\times 10^{-3}$. Dengan ambang satu, sebuah
peristiwa spontan akan terjadi pada kebanyakan jendelanya ($1 - e^{-1.25} = 0.71$)
dan gelapnya akan penuh kilatan; pada enam, tanda bahaya palsunya datang sekali
tiap lima ratus jendela. Ambangnya ditetapkan tempat derau
\omterm{def:b3:sensory-systems:retina}{sel batangnya}, bukan kepekaannya, yang menentukan.
\end{solution}

\begin{solution}{exo:b3:sensory-systems:11}
Penundaan terbesarnya $0.2/340 = \qty{590}{\micro s}$ (bunyi dari satu sisi).
Dekat garis tengahnya $\Delta t \approx (d/c)\sin\theta$, sehingga $\qty{10}{\micro
s}$ bersesuaian dengan $\sin\theta = 0.017$, yakni sekitar \qty{1}{\degree}. Di atas
\qty{4}{kHz} kalanya (\qty{250}{\micro s}) lebih pendek daripada yang dapat dikunci
neuronnya, dan sebuah penundaan beberapa ratus mikrodetik merentang
lebih daripada satu daur, yang membuat fasenya taksa; otaknya lalu memakai
selisih aras antara kedua telinganya dan sandi tempatnya.
\end{solution}

\begin{solution}{exo:b3:sensory-systems:12}
Jauh dari tepinya, $r = I(1 - 2w)$: $0.6$ di sisi gelapnya, $1.2$ di sisi
terangnya. Pada reseptor gelap yang terakhir ($I = 1$, tetangganya $1$ dan $2$): $r
= 1 - 0.2\times 3 = 0.4$; pada yang terang pertama: $r = 2 - 0.6 = 1.4$.
Tangganya dilebih-lebihkan menjadi sebuah lembah dan sebuah puncak --- yakni pita
Mach. \omterm{def:b3:sensory-systems:retina}{Retinanya} memperoleh kontras dan tepi, serta sebuah jangkauan dinamis yang lebih kecil
untuk ditularkan; ia kehilangan kesetiaan pada aras mutlak, lalu menghasilkan
ilusi kecerahan.
\end{solution}

\begin{solution}{pb:b3:sensory-systems:1}
\textbf{1.} $0.01/3.6\times 10^{-19} = 2.8\times 10^{16}$ foton tiap
detik.
\textbf{2.} Luas manik matanya $\pi(3.5\times 10^{-3})^{2} = 3.8\times 10^{-5}$
m$^{2}$. Pada \qty{1}{km}: pecahannya $3.8\times 10^{-5}/1.26\times 10^{7} =
3\times 10^{-12}$, yakni $8.5\times 10^{4}$ foton tiap detik. Pada \qty{10}{km}:
$850$ tiap detik.
\textbf{3.} Dalam \qty{0.1}{s} dengan \qty{10}{\%} terserap: $850$ pada
\qty{1}{km}, $8.5$ pada \qty{10}{km} --- keduanya di atas enam; lilinnya
terlihat pada sepuluh kilometer (hampir selalu).
\textbf{4.} Enam yang terserap tiap jendela memerlukan $600$ foton tiap detik ke dalam
manik matanya: $4\pi d^{2} = 3.8\times 10^{-5}\times 2.8\times 10^{16}/600
= 1.8\times 10^{9}$ m$^{2}$, $d \approx \qty{12}{km}$. Di sana, $P(k \ge
6 \mid a = 6) = 0.55$: lilinnya terlihat pada sekitar separuh jendelanya.
\textbf{5.} $e^{-6} = 0.0025$. Pada ambangnya cacahnya berayun dari
jendela ke jendela --- $6 \pm 2.4$ --- sehingga sumbernya seolah datang dan pergi:
yakni kerlipan sebuah bintang samar (yang padanya atmosfernya menambahkan
kerlipannya sendiri).
\textbf{6.} Ayunan latarnya, $\sqrt{10^{4}} = 100$ foton
tiap jendela, membanjiri $8.5$ milik lilinnya: pendeteksiannya menuntut isyaratnya
melampaui derau latarnya, bukan ambang mata yang gelap.
\textbf{7.} $\qty{1}{pA}/200 = \qty{5}{fA}$ tiap saluran; $5\times
10^{-15}/1.6\times 10^{-19} = 3\times 10^{4}$ ion tiap detik.
\textbf{8.} Sepertiga puluh; sekitar tiga puluh foton serentak menutup tiap
salurannya lalu menjenuhkan \omterm{def:b3:sensory-systems:retina}{sel batangnya}.
\textbf{9.} $10^{-12}\times 0.2 = 2\times 10^{-13}$ C, yakni $1.3\times 10^{6}$
kation: satu foton mengendalikan lebih daripada sejuta muatan.
\textbf{10.} \omterm{def:b3:sensory-systems:retina}{Sel batangnya} menjenuh dalam hitungan milidetik, tiap salurannya tertutup,
dan \omterm{def:b3:sensory-systems:retina}{rodopsinnya} terlunturkan lebih cepat daripada dipulihkan; \omterm{def:b3:sensory-systems:retina}{sel kerucut},
dengan bati lebih rendah dan pematian lebih cepat, terus menanggapi lalu menggeser
jangkauannya melintasi enam dasawarsa cahaya.
\textbf{11.} Bati yang tinggi memerlukan integrasi yang lama (tiap tahapnya memakan waktu dan
tanggapannya diperpanjang untuk menimbun), sehingga \omterm{def:b3:sensory-systems:retina}{sel batangnya} lambat dan
peka; pelipatgandaan \omterm{def:b3:sensory-systems:retina}{sel kerucut} yang lebih kecil dan pematiannya yang lebih cepat memberi sebuah
tanggapan dalam \qty{20}{ms}, cukup cepat bagi gerakan, dengan ongkos
memerlukan seratus foton.
\textbf{12.} $10^{8}/40 = 2.5\times 10^{6}$ isomerisasi spontan tiap
detik di seluruh \omterm{def:b3:sensory-systems:retina}{retinanya} --- tetapi hanya $1.25$ tiap jendela pada sembarang petak
$500$ \omterm{def:b3:sensory-systems:retina}{sel batang}, jauh di bawah ambang enamnya, sehingga deraunya ditolak
petak demi petak; ambangnya dan pengumpulannya ada tepat untuk hal ini.
\textbf{13.} \qty{0.1}{W/m^2}; pada \qty{5.5e-5}{m^2}, \qty{5.5e-6}{W}.
\textbf{14.} $5.5\times 10^{-6}\times 7200 = \qty{0.04}{J}$ --- yakni sekitar empat
ratus butir hujan sepanjang dua jam.
\textbf{15.} $10^{-12}\times 5.5\times 10^{-5} = 5.5\times 10^{-17}$ W;
dalam \qty{10}{ms}, $5.5\times 10^{-19}$ J --- yakni sekitar tenaga satu atau
dua foton tampak.
\textbf{16.} $0.3/5000 = 6\times 10^{-5}$ rad, yakni $0.003^{\circ}$;
lenturannya sekitar tiga atom hidrogen.
\textbf{17.} $10^{(110 - 85)/10} = 316$: dua jam pada \qty{110}{dB}
menghantarkan tenaga $630$ jam pada \qty{85}{dB} --- yakni hampir sebulan
hari kerja.
\textbf{18.} Sekitar $120$ langkah yang baru terasa, satu tiap \omterm{def:b3:sensory-systems:cochlea}{desibel}.
\textbf{19.} $2^{400} = 10^{400\log_{10}2} = 10^{120}$.
\textbf{20.} $\ln\binom{400}{20} \approx 400\,H(0.05) \approx 80$, dikurangi sebuah
koreksi Stirling sekitar $2$: yakni sekitar $10^{34}$ pola.
\textbf{21.} Cacahlah reseptor pada satu pola tetapi bukan pada yang lain: $5 +
5 = 10$ dari $40$ (yakni sebuah jarak Hamming); pola yang bertindih tiga
perempat menggerakkan \omterm{def:b3:sensory-systems:chemical}{glomerulus} yang nyaris sama lalu dibaca sebagai bau yang nyaris
sama.
\textbf{22.} $10^{12}/10^{34} = 10^{-22}$ dari polanya. Sandinya bukanlah
batasnya: reseptornya berderau, banyak yang ditala serupa sehingga
keaktifannya berkorelasi, molekul nyata menghasilkan hanya sebagian kecil
polanya, dan pembedaan serta ingatan otaknya atas pola itu terbatas.
\textbf{23.} $\ln\binom{1100}{20} \approx 1100\,H(0.018) = 1100\times
0.091 \approx 100$, yakni sekitar $10^{43}$ --- yakni sekitar $10^{9}$ kali angka
manusia bagi pola dua puluh reseptor (dan jauh lebih banyak bila mencitnya memakai
lebih banyak reseptor tiap bau).
\textbf{24.} Penciuman umumnya nyaris tidak berubah --- sandinya berlebihan dan
reseptor lain mengusung tiap zat bau --- tetapi sebuah zat bau yang bergantung pada
reseptor itu pada kepekatan rendah menjadi tak terdeteksi atau berubah
wataknya: yakni sebuah anosmia khas, seperti pada banyak orang yang tidak dapat
membaui androstenon.
\textbf{25.} Lilinnya terlihat sampai sekitar \qty{12}{km}; bunyi ambangnya
menghantarkan $5.5\times 10^{-19}$ J ke gendang telinganya dalam \qty{10}{ms}, yakni satu atau
dua foton; sebuah hidung manusia dapat membentuk sekitar $10^{34}$ pola dua puluh reseptor.
\end{solution}
